\section{Market handles}\label{sec:mp-market}
The market data the pricers and the models read: curves, volatilities, the equity forward. Every quote is an input
of the prices, and the sensitivities of a price are its derivatives w.r.t.\ these quotes.

% ─────────────────────────────────────────────────────────────────────────────
\subsection{Interest-rate curves}\label{sec:mp-curves}
A curve holds continuously compounded zero rates $r_i$ on pillars $T_i$ (Act/365), linear in between and flat outside the
pillars (definitions of \citealp[ch.~1]{brigo2006}):
\[
  P(0,T) = e^{-r(T)\,T},\qquad r(T) = r_i + \frac{r_{i+1}-r_i}{T_{i+1}-T_i}\,(T - T_i)\quad (T_i \le T < T_{i+1}).
\]
\begin{itemize}[nosep]
  \item zero rate $r(T) = -\ln P(0,T)/T$, with $r(0) = r_1$;
  \item simple forward $F(0;T_1,T_2) = \big(P(0,T_1)/P(0,T_2) - 1\big)/\tau$, $\tau = T_2 - T_1$ in the curve day count;
  \item instantaneous forward $f(0,T) = -\partial_T \ln P(0,T) = r(T) + T\,r'(T)$: at a pillar the right derivative,
        $f = r_1$ before the first pillar and $r_n$ after the last.
\end{itemize}
Each currency has an overnight curve (ESTR, SONIA, SOFR), the discount curve of everything in that currency; an Ibor
index has its own projection curve (EUR 6M) for its forward fixings.

\paragraph{Bootstrap.} The instruments are built from the quotes (deposits and OIS for the overnight curve; FRAs and
swaps for EUR 6M), with one pillar at the end date of each. The pillar rates are solved by Newton's method until every
instrument has zero NPV, each priced as in \S\ref{sec:mp-linear}; the Jacobian $\partial \mathrm{NPV}_i/\partial r_j$ of the final step links
the pillar sensitivities to the quote sensitivities.

% ─────────────────────────────────────────────────────────────────────────────
\subsection{Interest-rate volatilities}\label{sec:mp-ir-vols}
\begin{itemize}[nosep]
  \item \emph{Caplet vols}: normal (or lognormal) vols on fixing expiry $\times$ absolute strike nodes, bilinear, flat
        outside.
  \item \emph{Swaption cube}: vols on option expiry $\times$ swap tenor $\times$ strike offset $K - F$ from the forward swap rate,
        trilinear, flat outside.
  \item \emph{Caplet stripping}: from flat cap vols (one vol for all the caplets of a cap) to caplet vols, cap by cap in
        increasing maturity, each new caplet vol solving the price of the next cap; the inverse gives the flat vols of a
        caplet surface.
\end{itemize}

% ─────────────────────────────────────────────────────────────────────────────
\subsection{Equity spot, forward and dividends}\label{sec:mp-forward}
An equity has a spot $S_0$, a repo $q_r$, either cash dividends $D_i$ with ex-dates $t_i$ and payment dates $p_i$ or a
continuous yield $q$, and the discount curve $P(0,\cdot)$ of its currency. The forward, the same for every pricer and
model:
\begin{align*}
  \text{cash dividends:}\quad & F(T) = \frac{\big(S_0 - \mathrm{PV}_{\mathrm{div}}(T)\big)\,e^{-q_r T}}{P(0,T)},\qquad
    \mathrm{PV}_{\mathrm{div}}(T) = \sum_{t_i \le T} D_i\,P(0,p_i); \\
  \text{yield:}\quad & F(T) = \frac{S_0\,e^{-(q + q_r)T}}{P(0,T)}.
\end{align*}
A dividend is in the horizon by its ex-date (the stock drops then) and discounted from its payment date (the ex-date
if none is given). The equivalent dividend yield to $T$ is $q_{\mathrm{eff}}(T) = -\ln\big(1 - \mathrm{PV}_{\mathrm{div}}(T)/S_0\big)/T$. Cash
dividends are used when the schedule is given (the yield is then ignored), the yield otherwise.

% ─────────────────────────────────────────────────────────────────────────────
\subsection{Equity volatilities}\label{sec:mp-vol-surface}
Black implied vols $\sigma_{ij}$ on maturities $T_i$ (years) $\times$ absolute strikes $K_j$. For the closed-form pricers the vol at
$(K, T)$ is bilinear in $(K, T)$ and flat outside the grid, so it depends on the four surrounding nodes. The local vol model
builds its own representation from the same nodes, in total variance and forward log-moneyness (\S\ref{sec:mp-lv}).
